% 附录模板
% 使用方法：在 main.tex 的 \end{document} 前添加 % 附录模板
% 使用方法：在 main.tex 的 \end{document} 前添加 % 附录模板
% 使用方法：在 main.tex 的 \end{document} 前添加 % 附录模板
% 使用方法：在 main.tex 的 \end{document} 前添加 \input{appendix}

\appendix
\section{代码附录}

以下为本文的核心代码实现。

\subsection{问题一代码}

\begin{lstlisting}[language=Python, caption={问题一求解代码}]
% 在此粘贴问题一的代码
\end{lstlisting}

\subsection{问题二代码}

\begin{lstlisting}[language=Python, caption={问题二求解代码}]
% 在此粘贴问题二的代码
\end{lstlisting}

\subsection{问题三代码}

\begin{lstlisting}[language=Python, caption={问题三求解代码}]
% 在此粘贴问题三的代码
\end{lstlisting}

\section{数据附录}

以下为本文使用的数据样本。

\begin{table}[H]
\centering
\caption{数据样本}
\label{tab:data_sample}
\begin{tabular}{ccc}
\hline
列1 & 列2 & 列3 \\
\hline
数据1 & 数据2 & 数据3 \\
数据4 & 数据5 & 数据6 \\
\hline
\end{tabular}
\end{table}

\section{补充结果}

以下为补充实验结果。

\begin{figure}[H]
\centering
% \includegraphics[width=0.8\textwidth]{../figures/supplementary.pdf}
\caption{补充结果}
\label{fig:supplementary}
\end{figure}


\appendix
\section{代码附录}

以下为本文的核心代码实现。

\subsection{问题一代码}

\begin{lstlisting}[language=Python, caption={问题一求解代码}]
% 在此粘贴问题一的代码
\end{lstlisting}

\subsection{问题二代码}

\begin{lstlisting}[language=Python, caption={问题二求解代码}]
% 在此粘贴问题二的代码
\end{lstlisting}

\subsection{问题三代码}

\begin{lstlisting}[language=Python, caption={问题三求解代码}]
% 在此粘贴问题三的代码
\end{lstlisting}

\section{数据附录}

以下为本文使用的数据样本。

\begin{table}[H]
\centering
\caption{数据样本}
\label{tab:data_sample}
\begin{tabular}{ccc}
\hline
列1 & 列2 & 列3 \\
\hline
数据1 & 数据2 & 数据3 \\
数据4 & 数据5 & 数据6 \\
\hline
\end{tabular}
\end{table}

\section{补充结果}

以下为补充实验结果。

\begin{figure}[H]
\centering
% \includegraphics[width=0.8\textwidth]{../figures/supplementary.pdf}
\caption{补充结果}
\label{fig:supplementary}
\end{figure}


\appendix
\section{代码附录}

以下为本文的核心代码实现。

\subsection{问题一代码}

\begin{lstlisting}[language=Python, caption={问题一求解代码}]
% 在此粘贴问题一的代码
\end{lstlisting}

\subsection{问题二代码}

\begin{lstlisting}[language=Python, caption={问题二求解代码}]
% 在此粘贴问题二的代码
\end{lstlisting}

\subsection{问题三代码}

\begin{lstlisting}[language=Python, caption={问题三求解代码}]
% 在此粘贴问题三的代码
\end{lstlisting}

\section{数据附录}

以下为本文使用的数据样本。

\begin{table}[H]
\centering
\caption{数据样本}
\label{tab:data_sample}
\begin{tabular}{ccc}
\hline
列1 & 列2 & 列3 \\
\hline
数据1 & 数据2 & 数据3 \\
数据4 & 数据5 & 数据6 \\
\hline
\end{tabular}
\end{table}

\section{补充结果}

以下为补充实验结果。

\begin{figure}[H]
\centering
% \includegraphics[width=0.8\textwidth]{../figures/supplementary.pdf}
\caption{补充结果}
\label{fig:supplementary}
\end{figure}


\appendix
\section{代码附录}

以下为本文的核心代码实现。

\subsection{问题一代码}

\begin{lstlisting}[language=Python, caption={问题一求解代码}]
% 在此粘贴问题一的代码
\end{lstlisting}

\subsection{问题二代码}

\begin{lstlisting}[language=Python, caption={问题二求解代码}]
% 在此粘贴问题二的代码
\end{lstlisting}

\subsection{问题三代码}

\begin{lstlisting}[language=Python, caption={问题三求解代码}]
% 在此粘贴问题三的代码
\end{lstlisting}

\section{数据附录}

以下为本文使用的数据样本。

\begin{table}[H]
\centering
\caption{数据样本}
\label{tab:data_sample}
\begin{tabular}{ccc}
\hline
列1 & 列2 & 列3 \\
\hline
数据1 & 数据2 & 数据3 \\
数据4 & 数据5 & 数据6 \\
\hline
\end{tabular}
\end{table}

\section{补充结果}

以下为补充实验结果。

\begin{figure}[H]
\centering
% \includegraphics[width=0.8\textwidth]{../figures/supplementary.pdf}
\caption{补充结果}
\label{fig:supplementary}
\end{figure}
